% Generated by tables/make_robustness.py -- do not edit.
\begin{table}\caption{The estimate of $\theta$ across sub-samples}\label{tab:subsample}
\begin{singlespace}\begin{footnotesize}\begin{center}
\resizebox{\textwidth}{!}{\begin{tabular}{lcccccc}
\thickhline
Sample & \multicolumn{2}{c}{$G/Y$ weakly exogenous} & \multicolumn{2}{c}{LR $\alpha_{G/Y}=0$} & $\alpha$ free & $Q$ \\
 & $\hat\theta$ & $\hat\alpha_{\text{GAP}}$ & stat. & $p$ & $\hat\theta$ & $p$ \\
\thickhline
2003Q1--2019Q4 & 0.386 (0.030) & $-$0.109 (0.031) & 1.55 & 0.213 & 0.372 (0.032) & 0.66 \\
2005Q1--2019Q4 & 0.412 (0.037) & $-$0.099 (0.031) & 0.15 & 0.701 & 0.403 (0.043) & 0.59 \\
2007Q1--2019Q4 & 0.423 (0.053) & $-$0.079 (0.029) & 0.21 & 0.647 & 0.396 (0.079) & 0.74 \\
2010Q1--2019Q4 & 0.379 (0.093) & $-$0.076 (0.033) & 6.16 & 0.013 & 0.057 (0.144) & 0.30 \\
\thickhline\end{tabular}}\end{center}\end{footnotesize}
\begin{footnotesize}\textbf{Notes:} bivariate $(G/Y\mid\text{GAP})$, $p=4$, case 2, rank 1, on the same univariate models; the sample of the system starts at the date shown. LR: weak exogeneity of $G/Y$ (1 df). $Q$: Hosking's portmanteau of the adopted model.
\end{footnotesize}\end{singlespace}\end{table}
